\section{分式环}

记号约定:$A$是交换环,$B$是环,$S\subseteq A$,$A_{S}$是$A$相对$S$的分式幺半群(只配备乘法),$\pi:A\twoheadrightarrow A_{S}$是典范映射.

\begin{proposition}{分式环的构造}{construction-of-fraction-ring}
	$A_{S}$上仅有一种加法运算,使得配备该运算后$A_{S}$是交换环,$\pi\in\Hom_{\mathsf{Ring}}\left(A,A_{S}\right)$.
\end{proposition}

\begin{definition}{局部化,分式环}{}
	命题(\cref{prop:construction-of-fraction-ring})中构造的环称为$A$关于$S$的局部化(或$A$关于$S$的分式环),记作$A\left[S^{-1}\right]$.
	
	当$S$是$A$中所有可消去元组成的集合时,称$A\left[S^{-1}\right]$为$A$的全分式环(或$A$的完全局部化).
\end{definition}

\begin{remark}{记号说明}{}
	$S$是$A$中所有可消去元组成的集合时$\pi$是双射,因此在同构意义下$A$是$A\left[S^{-1}\right]$的子环.
\end{remark}

\begin{theorem}{局部化的泛性质}{}
	若$f\in\Hom_{\mathsf{Ring}}\left(A,B\right)$且$f\left(S\right)\subseteq B^{\times}$,则下图交换.
	\begin{center}
		\begin{tikzcd}
			A && {A\left[S^{-1}\right]} \\
			& B
			\arrow["\pi", two heads, from=1-1, to=1-3]
			\arrow["f"', from=1-1, to=2-2]
			\arrow["{\exists!\bar{f}}", dashed, from=1-3, to=2-2]
			\end{tikzcd}
	\end{center}
\end{theorem}


























